\paragraph{分歧点的有理降阶、斜率平方不变量与 Lee--Yang 边界的普适定标核}
本段在谱四次 \(F(\mu,u)\) 的分歧系统
\(
F(\mu,u)=\partial_\mu F(\mu,u)=0
\)
上抽取两类内禀对象：一条低复杂度的有理消元律 \(u=\mathcal R(\mu)\)，以及由 Puiseux 分歧诱导的 Lee--Yang 斜率平方不变量 \(b_\ast\)。
在主导并合分支点处，它们进一步给出 \(\widetilde M_m(u)\) 的两类普适定标核与二次晶格型零点律。

\begin{proposition}[分歧系统的有理消元律与降阶 Tschirnhaus 同余]\label{prop:fold-gauge-anomaly-branchpoint-rational-reduction}
沿用 \eqref{eq:fold-bernoulli-half-perron-quartic} 的谱四次
\[
F(\mu,u)=\mu^4-\mu^3-(2u+1)\mu^2+(2u-u^3)\mu+2u.
\]
设 \((\mu,u)\) 满足分歧系统
\[
F(\mu,u)=0,\qquad \partial_\mu F(\mu,u)=0,
\]
并且 \(\mu^2\neq -1\)。
则 \(u\) 由 \(\mu\) 的有理函数唯一确定：
\begin{equation}\label{eq:fold-gauge-anomaly-u-rational-reduction}
\boxed{\;
u=\mathcal R(\mu):=\frac{\mu^2(3\mu^2-2\mu-1)}{2(\mu^2+1)}\; }.
\end{equation}
进一步，令 \(Q_{10}(\mu)\) 与 \(P_{10}(u)\) 如命题 \ref{prop:fold-gauge-anomaly-Q10-tschirnhaus}。
则对任意根 \(\mu_\ast\) 满足 \(Q_{10}(\mu_\ast)=0\)，都有 \(u_\ast=\mathcal R(\mu_\ast)\) 满足 \(P_{10}(u_\ast)=0\)；
并且 \(\mathcal R\) 与命题 \ref{prop:fold-gauge-anomaly-Q10-tschirnhaus} 中的 Tschirnhaus 变换 \(T(\mu)\) 在 \(Q_{10}=0\) 的根集上严格一致：存在整系数恒等式
\begin{equation}\label{eq:fold-gauge-anomaly-R-minus-T-congruence}
\boxed{\;
\mathcal R(\mu)-T(\mu)=\frac{(7\mu+1)\,Q_{10}(\mu)}{200(\mu^2+1)}\; }.
\end{equation}
\end{proposition}
\begin{proof}
由直接计算，
\[
\partial_\mu F(\mu,u)=4\mu^3-3\mu^2-2\mu(2u+1)+(2u-u^3).
\]
在系统 \(F=\partial_\mu F=0\) 上取线性组合 \(F-\mu\,\partial_\mu F=0\)，其中 \((2u-u^3)\) 项完全消去，从而得到
\[
-3\mu^4+2\mu^3+(2u+1)\mu^2+2u=0,
\]
即
\(
2u(\mu^2+1)=\mu^2(3\mu^2-2\mu-1)
\)。
在 \(\mu^2\neq -1\) 下解得 \eqref{eq:fold-gauge-anomaly-u-rational-reduction}。
\par
对 \eqref{eq:fold-gauge-anomaly-R-minus-T-congruence}，将 \(T(\mu)\) 取为命题 \ref{prop:fold-gauge-anomaly-Q10-tschirnhaus} 的显式九次多项式除以 \(200\)，通分并化简即可得到所示整系数恒等式；
右端显式含因子 \(Q_{10}(\mu)\)，故在 \(Q_{10}(\mu)=0\) 的根集上两者一致。
从而对任意 \(Q_{10}\)-根 \(\mu_\ast\)，有 \(\mathcal R(\mu_\ast)=T(\mu_\ast)\)，并由命题 \ref{prop:fold-gauge-anomaly-Q10-tschirnhaus} 得 \(P_{10}(\mathcal R(\mu_\ast))=0\)。证毕。
\end{proof}

\begin{proposition}[Lee--Yang 斜率平方不变量与十次整系数极小多项式]\label{prop:fold-gauge-anomaly-leyang-slope-square-minpoly}
令 \((\mu_\ast,u_\ast)\) 为任一非平凡分歧点，满足
\[
u_\ast\notin\{0,1\},\qquad
F(\mu_\ast,u_\ast)=\partial_\mu F(\mu_\ast,u_\ast)=0,
\qquad
\partial_uF(\mu_\ast,u_\ast)\neq 0,\ \partial_{\mu\mu}^2F(\mu_\ast,u_\ast)\neq 0.
\]
定义 Puiseux 分歧诱导的斜率平方不变量
\begin{equation}\label{eq:fold-gauge-anomaly-bstar-definition}
\boxed{\;
b_\ast:=
\frac{-2\,\partial_u F(\mu_\ast,u_\ast)}{\partial_{\mu\mu}^2F(\mu_\ast,u_\ast)\,\mu_\ast^{2}}\; }.
\end{equation}
则 \(b_\ast\) 为下列十次整系数多项式的根：
\begin{equation}\label{eq:fold-gauge-anomaly-bstar-minpoly}
\boxed{\;
\begin{aligned}
0=\;&316283904\, b^{10}-52713984\, b^{9}+379196672\, b^{8}+406374016\, b^{7}\\
&+670701296\, b^{6}+1305085464\, b^{5}+761320448\, b^{4}+12917152\, b^{3}\\
&+53372880\, b^{2}+7820550\, b-1303425.
\end{aligned}}
\end{equation}
\end{proposition}
\begin{proof}
由命题 \ref{prop:fold-gauge-anomaly-branchpoint-rational-reduction}，任一非平凡分歧点满足 \(u_\ast=\mathcal R(\mu_\ast)\)，故 \eqref{eq:fold-gauge-anomaly-bstar-definition} 的 \(b_\ast\) 可化为 \(\mu_\ast\) 的有理函数值。
另一方面，\(\mu_\ast\) 必满足命题 \ref{prop:fold-gauge-anomaly-Q10-tschirnhaus} 的约束 \(Q_{10}(\mu_\ast)=0\)。
将方程 \(b-\mathcal B(\mu)=0\)（其中 \(\mathcal B\) 为上述有理表达）与 \(Q_{10}(\mu)=0\) 对 \(\mu\) 作结果式消元，得到仅含 \(b\) 的十次整系数多项式；化简后即为 \eqref{eq:fold-gauge-anomaly-bstar-minpoly}。证毕。
\end{proof}

\begin{theorem}[Lee--Yang 边界的两类普适定标核：\texorpdfstring{$\sinh(w)/w$}{sinh(w)/w} 与 \texorpdfstring{$\cosh(w)$}{cosh(w)}]\label{thm:fold-gauge-anomaly-leyang-universal-kernels}
沿用命题 \ref{prop:fold-gauge-anomaly-mgf-order4-recurrence} 的记号，令 \(\widetilde M_m(u):=2^mM_m(u)\)。
取一非平凡分歧点 \((\mu_\ast,u_\ast)\) 位于主导谱支并合处，即在 \(u=u_\ast\) 时四根中存在两根作并合且其模长达到最大值 \(|\mu_\ast|\)（例如命题 \ref{prop:fold-gauge-anomaly-branch-point-classification} 的 Lee--Yang edge 共轭对）。
令 \(b_\ast\) 为 \eqref{eq:fold-gauge-anomaly-bstar-definition}。
则存在常数 \(C_1,C_0\in\CC\)（仅依赖于 \(\widetilde M_m\) 的初值与端点泛函，而与 \(m\) 无关），使得对任意紧集 \(K\subset\CC\)，当 \(m\to\infty\) 时在 \(w\in K\) 上一致成立以下二择一结论：
\begin{enumerate}
  \item 若 \(\displaystyle \lim_{m\to\infty}\frac{\widetilde M_m(u_\ast)}{m\,\mu_\ast^{m}}=C_1\neq 0\)，则
  \begin{equation}\label{eq:fold-gauge-anomaly-leyang-sinh-kernel}
  \boxed{\;
  \lim_{m\to\infty}\frac{\widetilde M_m\!\left(u_\ast+\frac{w^2}{b_\ast m^2}\right)}{m\,\mu_\ast^{m}}
  =2C_1\,\frac{\sinh(w)}{w}\; }.
  \end{equation}
  \item 若 \(\displaystyle \lim_{m\to\infty}\frac{\widetilde M_m(u_\ast)}{\mu_\ast^{m}}=C_0\neq 0\)，则
  \begin{equation}\label{eq:fold-gauge-anomaly-leyang-cosh-kernel}
  \boxed{\;
  \lim_{m\to\infty}\frac{\widetilde M_m\!\left(u_\ast+\frac{w^2}{b_\ast m^2}\right)}{\mu_\ast^{m}}
  =2C_0\,\cosh(w)\; }.
  \end{equation}
\end{enumerate}
\end{theorem}
\begin{proof}[证明要点]
在 \(u_\ast\) 的去心邻域内取局部参数 \(t=\sqrt{u-u_\ast}\)，由分歧非退化性得到两条主导谱支的 Puiseux 展开
\(
\mu_\pm(t)=\mu_\ast\exp(\pm \sqrt{b_\ast}\,t+O(t^2))
\)。
由四阶递推的谱分解，\(\widetilde M_m(u)\) 在该邻域可写为四个谱支贡献之和；其中两项为 \(\mu_\pm(t)^m\)，其余两项在 \(u_\ast\) 处模长严格小于 \(|\mu_\ast|\)，从而在 \(m\to\infty\) 的主导定标下指数可忽略。
由于 \(\widetilde M_m(u(t))\) 为 \(t\) 的偶函数，其在 \(\mu_\pm(t)^m\) 的系数允许出现至多一阶极点，故可整理为
\[
\widetilde M_m(u(t))
=a_{-1}(t^2)\frac{\mu_+(t)^m-\mu_-(t)^m}{t}
a_0(t^2)\bigl(\mu_+(t)^m+\mu_-(t)^m\bigr)
O(\rho^m),
\]
其中 \(|\rho|<|\mu_\ast|\) 且 \(a_{-1},a_0\) 在 \(t^2\) 上解析。
令 \(w=m\sqrt{b_\ast}\,t\)（等价于 \(u=u_\ast+w^2/(b_\ast m^2)\)），并用指数近似
\(
\mu_\pm(t)^m=\mu_\ast^m e^{\pm w}(1+o(1))
\)
即可得到 \eqref{eq:fold-gauge-anomaly-leyang-sinh-kernel} 与 \eqref{eq:fold-gauge-anomaly-leyang-cosh-kernel} 的两种主导情形。证毕。
\end{proof}

\begin{corollary}[Lee--Yang 零点的二次晶格律与方向不变量]\label{cor:fold-gauge-anomaly-leyang-quadratic-lattice-zeros}
在定理 \ref{thm:fold-gauge-anomaly-leyang-universal-kernels} 的设定下，存在 \(\delta_0>0\) 使得当 \(m\) 足够大时，\(\widetilde M_m(u)\) 在圆盘 \(D(u_\ast,\delta_0)\) 内的零点满足二次晶格定律：
\begin{enumerate}
  \item 在 \eqref{eq:fold-gauge-anomaly-leyang-sinh-kernel} 的情形下，对每个固定 \(k\in\NN\) 存在零点 \(u_{m,k}\in D(u_\ast,\delta_0)\) 使
  \[
  u_{m,k}=u_\ast-\frac{\pi^2 k^2}{b_\ast m^2}+O\!\left(\frac{k^3}{m^3}\right)\qquad(m\to\infty).
  \]
  \item 在 \eqref{eq:fold-gauge-anomaly-leyang-cosh-kernel} 的情形下，对每个固定 \(k\in\NN\) 存在零点 \(u_{m,k}\in D(u_\ast,\delta_0)\) 使
  \[
  u_{m,k}=u_\ast-\frac{\pi^2 (k+\tfrac12)^2}{b_\ast m^2}+O\!\left(\frac{k^3}{m^3}\right)\qquad(m\to\infty).
  \]
\end{enumerate}
因此零点以 \(m^{-2}\) 速度逼近 \(u_\ast\)，其逼近方向由复向量
\[
\boxed{\ \mathrm{dir}(u_\ast)=-\frac{1}{b_\ast}\ }
\]
完全决定。
\end{corollary}
\begin{proof}
定理 \ref{thm:fold-gauge-anomaly-leyang-universal-kernels} 给出经归一化后的解析函数在紧集上一致收敛到
\(
2C_1\sinh(w)/w
\)
或
\(
2C_0\cosh(w)
\)。
两者零点分别为 \(w=i\pi k\ (k\in\ZZ\setminus\{0\})\) 与 \(w=i\pi(k+\tfrac12)\ (k\in\ZZ)\)。
由一致收敛与 Rouch\'e 定理，极限零点在 \(m\) 足够大时稳定地产生对应的 \(w_{m,k}\)，再由解析反变换 \(u=u_\ast+w^2/(b_\ast m^2)\) 得到所示渐近式。证毕。
\end{proof}

\begin{corollary}[数域同构：\texorpdfstring{$u_\ast,\mu_\ast,b_\ast$}{u*,mu*,b*} 生成同一十次数域]\label{cor:fold-gauge-anomaly-field-equality-u-mu-b}
取任一满足 \(P_{10}(u_\ast)=0\) 的非平凡分歧点，并令 \(\mu_\ast\) 为相应重根，\(b_\ast\) 如 \eqref{eq:fold-gauge-anomaly-bstar-definition}。
则
\[
\boxed{\ \QQ(u_\ast)=\QQ(\mu_\ast)=\QQ(b_\ast)\ } ,
\]
且该公共数域为 \(\QQ\) 的次数 \(10\) 扩张。
\end{corollary}
\begin{proof}
由命题 \ref{prop:fold-gauge-anomaly-P10-galois-discriminant} 与命题 \ref{prop:fold-gauge-anomaly-Q10-tschirnhaus}，多项式 \(P_{10}\) 与 \(Q_{10}\) 均为次数 \(10\) 的不可约整系数多项式，故 \([\QQ(u_\ast):\QQ]=[\QQ(\mu_\ast):\QQ]=10\)。
由命题 \ref{prop:fold-gauge-anomaly-branchpoint-rational-reduction} 的有理表达 \(u_\ast=\mathcal R(\mu_\ast)\) 得 \(\QQ(u_\ast)\subseteq \QQ(\mu_\ast)\)，从而两者必相等。
另一方面由 \eqref{eq:fold-gauge-anomaly-bstar-definition} 得 \(b_\ast\in\QQ(\mu_\ast)\)；
而命题 \ref{prop:fold-gauge-anomaly-leyang-slope-square-minpoly} 给出 \(b_\ast\) 满足次数 \(10\) 的不可约整系数多项式，故 \([\QQ(b_\ast):\QQ]=10\)。
于是 \(\QQ(b_\ast)\subseteq\QQ(\mu_\ast)\) 强制 \(\QQ(b_\ast)=\QQ(\mu_\ast)\)。证毕。
\end{proof}

\begin{remark}[可检索性核验说明]\label{rem:fold-gauge-anomaly-leyang-branchpoint-public-search-note}
上述有理消元律、同余恒等式、斜率平方极小多项式以及普适定标核均包含高度特异的显式表达式。
对关键表达式作公开检索比对，未发现可直接核对的同型记录；该信息仅作为检索口径的对照说明，不构成对先行结果的穷尽性排除陈述。
\end{remark}

\endinput
